\subsection{RQ2: which node view transfers across projects}
\label{sec:rq_representation}

\emph{In short:} on projects never seen in training, a network that reads
identifier text from scratch performs close to chance; removing the project
vocabulary (by renaming identifiers or by using grammar kinds) lifts AUC by
roughly 0.1 to 0.17; and the choice among the vocabulary-free views matters
much less than the choice to leave the vocabulary out.

Everything except the node view is held fixed: graph, backbone (GGNN), depth,
width, readout, optimiser, class weighting and gradient budget.
Table~\ref{tab:rq2_representation} reports the BigVul results over 10 matched
folds, and Table~\ref{tab:direct} the direct paired contrasts between views.

% Generated by experiments/emit_tables.py -- DO NOT EDIT BY HAND.
% Every value is read from experiments/results/*.json.

\begin{table*}[t]
\centering
\caption{RQ2: node views on BigVul, GGNN, 10 matched project-level folds
(mean $\pm$ population std over folds). Differences are paired against the
lexical baseline: each $\Delta$ cell reads mean per-fold difference, 95\%
cluster-bootstrap interval over held-out projects, bootstrap $p$ / Wilcoxon
signed-rank $p$ over folds; \emph{n.s.} marks an interval containing zero.
\emph{AUC in size bins} is the AUC computed within five bins of function size
(0.5 = no information beyond length). AUPRC is reported without paired tests.}
\label{tab:rq2_representation}
\resizebox{\linewidth}{!}{%
\begin{tabular}{lrrlrr}
\toprule
Node view & Folds & AUC & $\Delta$AUC [95\% CI] $p_{\text{boot}}$ / $p_{\text{W}}$ & AUC in size bins & AUPRC \\
\midrule
Trivial (all-positive) & --- & 0.500 & --- & 0.500 & 0.071 \\
Function size only (rule) & 10 & 0.736 $\pm$ 0.065 & --- & 0.500 & 0.232 $\pm$ 0.087 \\
Lexical token & 10 & 0.506 $\pm$ 0.077 & --- & 0.467 & 0.085 $\pm$ 0.034 \\
Anonymised token & 10 & 0.708 $\pm$ 0.046 & $+0.202$ $[+0.149, +0.249]$ $<0.001$ / 0.002 & 0.569 & 0.167 $\pm$ 0.048 \\
Operation type $\phi_8$ & 10 & 0.649 $\pm$ 0.052 & $+0.143$ $[+0.092, +0.211]$ $<0.001$ / 0.010 & 0.497 & 0.117 $\pm$ 0.026 \\
Grammar kind $\kappa$ & 10 & 0.691 $\pm$ 0.037 & $+0.185$ $[+0.130, +0.236]$ $<0.001$ / 0.002 & 0.535 & 0.149 $\pm$ 0.027 \\
Operation $+$ kind ($\kappa{+}\phi$) & 10 & 0.687 $\pm$ 0.060 & $+0.181$ $[+0.132, +0.236]$ $<0.001$ / 0.002 & 0.536 & 0.158 $\pm$ 0.042 \\
Lexical $+$ kind & 10 & 0.592 $\pm$ 0.072 & $+0.086$ $[+0.038, +0.111]$ $<0.001$ / 0.002 & 0.519 & 0.105 $\pm$ 0.039 \\
Anonymised $+$ kind & 10 & 0.701 $\pm$ 0.059 & $+0.195$ $[+0.144, +0.244]$ $<0.001$ / 0.002 & 0.558 & 0.180 $\pm$ 0.059 \\
\bottomrule
\end{tabular}}
\end{table*}


\paragraph{Text read from scratch does not transfer.} With the lexical view
the GGNN reaches AUC \mBvGgnnLexAuc{} on BigVul and \mDvGgnnLexAuc{} on
DiverseVul, close to the $0.5$ of a random ranking. The cause is visible in the
mechanism: after the per-fold cut, tokens that occur only in held-out projects
map to \texttt{<UNK>}, and what remains of the vocabulary correlates with
project identity rather than with the weakness.

\paragraph{Vocabulary-free views transfer.} Replacing text with the grammar kind
raises AUC by \cBvGgnnKindVLexD{} (95\% CI \cBvGgnnKindVLexLo{} to
\cBvGgnnKindVLexHi{}; bootstrap $p$ \cBvGgnnKindVLexP{}; Wilcoxon $p$
\cBvGgnnKindVLexW{}; \cBvGgnnKindVLexWins{} folds won) on BigVul, and by
\cDvGgnnKindVLexD{} (\cDvGgnnKindVLexLo{} to \cDvGgnnKindVLexHi{}) on
DiverseVul. Renaming identifiers while keeping all other token text
(\emph{anonymised}) gives \cBvGgnnAnonVLexD{}
(\cBvGgnnAnonVLexLo{} to \cBvGgnnAnonVLexHi{}) on BigVul. The effect is large
relative to the fold-to-fold spread of the lexical view
(standard deviation \mBvGgnnLexAucStd{} AUC; Table~\ref{tab:rq2_representation}).

\paragraph{The grammar kind is not better than anonymised text.} This is the
conclusion that changed most since the first version of this work. Reviewing it
we added the baselines that a text-only comparison omits: anonymised tokens, and
each of text and anonymised tokens combined with the kind
(Section~\ref{sec:features}). The direct contrast of $\kappa$ against
anonymised$+\kappa$ is \cBvGgnnKindVAnonKindD{}
(\cBvGgnnKindVAnonKindLo{} to \cBvGgnnKindVAnonKindHi{}) on BigVul and
\cDvGgnnKindVAnonKindD{} (\cDvGgnnKindVAnonKindLo{} to
\cDvGgnnKindVAnonKindHi{}) on DiverseVul
(Table~\ref{tab:direct}). Identifier renaming, which is what sequence-based
detectors have long done as preprocessing, therefore recovers essentially the
whole benefit. What the grammar kind offers is the same effect without any
preprocessing choice, not a better effect.

\paragraph{Adding node type to text helps, but less than removing the text.}
The Devign-style combination of text with the node type improves on text alone
by \cBvGgnnLexKindVLexD{} (\cBvGgnnLexKindVLexLo{} to
\cBvGgnnLexKindVLexHi{}) on BigVul, so part of the benefit of structure is
independent of removing names. But $\kappa$ alone beats that combination by
\cBvGgnnKindVLexKindD{} (\cBvGgnnKindVLexKindLo{} to
\cBvGgnnKindVLexKindHi{}): while project-specific text is present, the network
keeps fitting it. The comparison in the first submission, which pitted
$\kappa$ against a text view that carried no node-type information at all, was
therefore partly ``structure labels versus none''; Table~\ref{tab:direct} now
separates the two.

\paragraph{The engineered operation taxonomy adds nothing to the kind.}
The direct paired contrast of $\kappa{+}\phi$ against $\kappa$ is
\cBvGgnnOpKindVKindD{} (\cBvGgnnOpKindVKindLo{} to \cBvGgnnOpKindVKindHi{};
Wilcoxon $p$ \cBvGgnnOpKindVKindW{}) on BigVul and \cDvGgnnOpKindVKindD{}
(\cDvGgnnOpKindVKindLo{} to \cDvGgnnOpKindVKindHi{}) on DiverseVul. The
earlier submission inferred ``no benefit'' from the two views' separate gaps
to the lexical baseline, which is not a test; here it is tested directly, and
the interval contains zero in both corpora. $\phi$ alone ($\phi_8$) is
\cBvGgnnOpVLexD{} above text and \cBvGgnnKindVOpD{} below $\kappa$
(\cBvGgnnKindVOpLo{} to \cBvGgnnKindVOpHi{}); the second difference is
estimated imprecisely and we do not claim it. Section~\ref{sec:rq_why}
explains why a null here is expected.

\paragraph{Capacity control.} The kind table has 364 rows and the taxonomy 9, so
one might suspect that $\kappa$ wins by capacity. We therefore also trained
$\kappa$ squeezed through a nine-dimensional bottleneck, the width of the
taxonomy table (\texttt{kind\_low}). Its AUC is \mBvGgnnKindLowAuc{}, and the
difference from the unconstrained kind embedding is \cBvGgnnKindLowVKindD{}
(\cBvGgnnKindLowVKindLo{} to \cBvGgnnKindLowVKindHi{}), consistent with no
capacity effect. (An earlier version of this paper reported a similar check
without any code or data behind it; it has been replaced by this experiment.)

\paragraph{Controlling for function size.} Vulnerable functions are longer
(Section~\ref{sec:corpora}), so a detector can gain AUC from length alone. A
rule that scores each function by its node count reaches AUC \numbvSizeAuc{} on
BigVul and \numdvSizeAuc{} on DiverseVul (second row of
Table~\ref{tab:rq2_representation}), higher than every GNN here. To ask what
the networks learn beyond length we also compute the AUC \emph{within} five
bins of function size (pooled held-out predictions, bin edges from the
held-out size quantiles); a score that carries no information beyond length
has within-bin AUC $0.5$. On BigVul the grammar kind's within-bin AUC is
\mBvGgnnKindSs{} against \mBvGgnnLexSs{} for text, a paired difference of
\cBvGgnnKindVLexSsD{} (\cBvGgnnKindVLexSsLo{} to \cBvGgnnKindVLexSsHi{}). The
gain over text is therefore only partly independent of length, and the
within-bin skill of every view is modest.

\paragraph{No single project drives the result.}
Because one repository can dominate a fold, we repeat the BigVul contrast with
that project's test predictions removed: $\kappa$ against text is
\cBvGgnnKindVLexNoAnd{} without \texttt{\numbigvulTopTestProject}, compared with
\cBvGgnnKindVLexD{} with it. Table~\ref{tab:perproject} gives the per-project
AUCs for the held-out projects with at least ten vulnerable functions:
the pattern of $\kappa$ over text is present in most of them, although each
single-project figure rests on few positives and is noisy.

% Generated by experiments/emit_tables.py -- DO NOT EDIT BY HAND.
% Every value is read from experiments/results/*.json.

\begin{table}[t]
\centering
\caption{Per-project AUC (BigVul, GGNN) for the held-out projects with at
least 10 vulnerable functions, pooled over the folds in which they are tested.
It shows that the pooled result is not carried by a single repository.}
\label{tab:perproject}
\resizebox{\linewidth}{!}{%
\begin{tabular}{lrrrrrr}
\toprule
Project & Functions & Vuln. & Lexical & Lex$+\kappa$ & $\kappa$ & $\kappa{+}\phi$ \\
\midrule
Android & 434 & 60 & 0.50 & --- & 0.64 & 0.55 \\
ImageMagick & 135 & 21 & 0.59 & --- & 0.67 & 0.66 \\
\bottomrule
\end{tabular}}
\end{table}

% Generated by experiments/emit_tables.py -- DO NOT EDIT BY HAND.
% Every value is read from experiments/results/*.json.

\begin{table*}[t]
\centering
\caption{Direct paired contrasts between node views (GGNN, 10 matched folds):
$\Delta$AUC of the first view minus the second, with 95\% cluster-bootstrap
interval, bootstrap $p$ / Wilcoxon $p$. These are the comparisons a reader
needs to judge whether one view is better than another; a difference between
two views' gaps to the lexical baseline is not a test of the difference.}
\label{tab:direct}
\resizebox{\linewidth}{!}{%
\begin{tabular}{lll}
\toprule
Contrast & BigVul & DiverseVul \\
\midrule
$\kappa{+}\phi$ vs $\kappa$ & $-0.004$ $[-0.017, +0.021]$ 0.920\,n.s. / 1.000 & $+0.008$ $[-0.005, +0.022]$ 0.232\,n.s. / 0.027 \\
$\kappa$ vs lexical$+\kappa$ & $+0.100$ $[+0.064, +0.156]$ $<0.001$ / 0.002 & $+0.060$ $[+0.030, +0.088]$ $<0.001$ / 0.002 \\
$\kappa$ vs anonymised$+\kappa$ & $-0.010$ $[-0.038, +0.014]$ 0.380\,n.s. / 0.625 & $-0.012$ $[-0.030, +0.002]$ 0.098\,n.s. / 0.322 \\
anonymised$+\kappa$ vs lexical$+\kappa$ & $+0.110$ $[+0.079, +0.166]$ $<0.001$ / 0.002 & $+0.071$ $[+0.043, +0.102]$ $<0.001$ / 0.004 \\
lexical$+\kappa$ vs lexical & $+0.086$ $[+0.038, +0.111]$ $<0.001$ / 0.002 & $+0.042$ $[+0.021, +0.061]$ $<0.001$ / 0.037 \\
anonymised vs lexical & $+0.202$ $[+0.149, +0.249]$ $<0.001$ / 0.002 & --- \\
$\kappa$ vs $\phi_8$ & $+0.042$ $[-0.004, +0.067]$ 0.072\,n.s. / 0.049 & --- \\
9-dim bottleneck $\kappa$ vs $\kappa$ & $+0.017$ $[+0.004, +0.040]$ 0.018 / 0.322 & --- \\
\bottomrule
\end{tabular}}
\end{table*}

